% TEMPLATE-MILC-v3.tex - plantilla maestra para todas las guias MILC v3
% Ver scripts/generadores/build-guia-milc.py para la lista completa de
% claves esperadas. El script valida que no queden placeholders sin
% reemplazar antes de escribir el archivo final.

\documentclass[11pt,letterpaper]{article}
\usepackage[margin=1.75cm]{geometry}
\usepackage[table]{xcolor}
\usepackage{fontspec}
\usepackage[spanish]{babel}
\usepackage{tikz}
\usepackage[most]{tcolorbox}
\usepackage{tabularx}
\usepackage{array}
\usepackage{fancyhdr}
\usepackage{titlesec}
\usepackage{ragged2e}
\usepackage{enumitem}
\usepackage{microtype}

\setmainfont{Helvetica Neue}
\setsansfont{Helvetica Neue}
\setlength{\parindent}{0pt}
\setlength{\parskip}{6pt}
\hyphenpenalty=200
\exhyphenpenalty=200
\pretolerance=400
\tolerance=2000
\setlength{\emergencystretch}{1em}
\renewcommand{\arraystretch}{1.25}
\setlist{leftmargin=5mm,itemsep=1mm,topsep=1mm}
\newcolumntype{Y}{>{\RaggedRight\arraybackslash}X}

\definecolor{milcMagenta}{HTML}{E6007E}
\definecolor{milcVino}{HTML}{5A0038}
\definecolor{milcTurquesa}{HTML}{0093A5}
\definecolor{milcVerde}{HTML}{58A923}
\definecolor{milcMostaza}{HTML}{E5B400}
\definecolor{milcOcre}{HTML}{B86600}
\definecolor{milcNegro}{HTML}{241816}
\definecolor{milcGris}{HTML}{F7F7F4}
\definecolor{milcLinea}{HTML}{E8E2DD}
\definecolor{milcGradePrimary}{HTML}{7C3AED}
\definecolor{milcGradeDark}{HTML}{4C1D95}
\definecolor{milcGradeSoft}{HTML}{EDE9FE}
\definecolor{milcGradeAccent}{HTML}{CFE7FF}
\definecolor{milcGradeText}{HTML}{FFFFFF}
\color{milcNegro}

\pagestyle{fancy}
\fancyhf{}
\lhead{\small Tecnología e Informática · Grado 9}
\rhead{\small Guía 16 · MILC v3}
\cfoot{\small\thepage}
\renewcommand{\headrulewidth}{0pt}

\newtcolorbox{titlebox}[1]{
  enhanced,colback=#1,colframe=#1,arc=7mm,boxrule=0pt,
  left=7mm,right=7mm,top=3mm,bottom=3mm,
  before skip=8pt,after skip=8pt,
  fontupper=\sffamily\bfseries\Large\color{white}
}

\newtcolorbox{softbox}[3]{
  enhanced,colback=#2,colframe=#1,borderline west={4pt}{0pt}{#1},
  arc=2mm,boxrule=.4pt,left=4mm,right=4mm,top=3mm,bottom=3mm,
  before skip=6pt,after skip=8pt,title={#3},coltitle=white,
  colbacktitle=#1,fonttitle=\sffamily\bfseries,fontupper=\RaggedRight,
  attach boxed title to top left={xshift=3mm,yshift=-2mm},
  boxed title style={arc=2mm, boxrule=0pt}
}

\newtcolorbox{drawbox}[2]{
  enhanced,colback=white,colframe=#1,arc=4mm,boxrule=1pt,
  left=5mm,right=5mm,top=4mm,bottom=4mm,before skip=8pt,after skip=8pt,
  title={#2},coltitle=white,colbacktitle=#1,fonttitle=\sffamily\bfseries,
  fontupper=\RaggedRight,
  attach boxed title to top left={xshift=4mm,yshift=-2mm},
  boxed title style={arc=3mm, boxrule=0pt}
}

\newtcolorbox{infoband}[1]{
  enhanced,colback=#1!8,colframe=#1!50,arc=2mm,boxrule=.5pt,
  left=4mm,right=4mm,top=2mm,bottom=2mm,before skip=4pt,after skip=4pt,
  fontupper=\small\RaggedRight
}

% Puente narrativo entre secciones (contrato editorial Conéctate).
% Da continuidad: "Acabas de X. Ahora vas a Y porque Z."
% Visualmente: banda izquierda en color del grado, texto en itálica pequeña.
\newtcolorbox{puentebox}{
  enhanced,colback=milcGradeSoft,colframe=milcGradeDark,
  borderline west={3pt}{0pt}{milcGradeDark},
  arc=0mm,boxrule=0pt,left=5mm,right=4mm,top=2.5mm,bottom=2.5mm,
  before skip=4pt,after skip=8pt,
  fontupper=\itshape\small\color{milcNegro}\RaggedRight
}
\newcommand{\puente}[1]{%
  \begin{puentebox}%
    \textcolor{milcGradeDark}{\textbf{\textrightarrow}}\hspace{2mm}#1%
  \end{puentebox}%
}

\newcommand{\linead}[1]{\noindent\color{milcLinea}\rule{#1}{0.5pt}\color{milcNegro}}
\newcommand{\respuesta}[1]{\par\vspace{2mm}\foreach \n in {1,...,#1}{\noindent\color{milcLinea}\rule{\linewidth}{0.5pt}\par\vspace{5mm}}\color{milcNegro}}
\newcommand{\checkbox}{\textcolor{milcGradeDark}{\fbox{\phantom{x}}}\hspace{5pt}}

\newcommand{\phasecard}[4]{
\begin{tcolorbox}[enhanced,colback=#3,colframe=#2,arc=3mm,boxrule=.5pt,height=3.05cm,valign=center,halign=center,left=1.5mm,right=1.5mm,top=2mm,bottom=2mm]
{\sffamily\bfseries\fontsize{22}{24}\selectfont\textcolor{#2}{#1}}\par
{\sffamily\bfseries\small #4}
\end{tcolorbox}
}

\begin{document}
\thispagestyle{empty}
\begin{tikzpicture}[remember picture,overlay]
  \fill[white] (current page.south west) rectangle (current page.north east);
  \path[fill=milcGradePrimary,rounded corners=16mm]
    ([xshift=.55cm,yshift=-.55cm]current page.north west) rectangle
    ([xshift=-.55cm,yshift=3.65cm]current page.south east);

  \node[anchor=north west,text=milcGradeText] at ([xshift=2.0cm,yshift=-1.85cm]current page.north west)
    {\sffamily\bfseries\fontsize{14}{16}\selectfont GRADO};
  \node[anchor=north west,text=milcGradeText,opacity=.82] at ([xshift=2.0cm,yshift=-2.75cm]current page.north west)
    {\sffamily\bfseries\fontsize{9}{11}\selectfont SERIE GUÍAS · TECNOLOGÍA E INFORMÁTICA};

  \begin{scope}[shift={([xshift=-3.05cm,yshift=-2.55cm]current page.north east)}]
    \fill[white,opacity=.80] (-.62,-.52) rectangle (.62,.52);
    \draw[milcGradeText,line width=1pt] (-.62,-.52) rectangle (.62,.52);
    \fill[milcGradeDark] (-.42,-.38) rectangle (-.22,.06);
    \fill[milcGradeAccent] (-.10,-.38) rectangle (.10,.24);
    \fill[milcGradePrimary!60!black] (.22,-.38) rectangle (.42,.38);
  \end{scope}

  \node[anchor=west,text=milcGradeText] at ([xshift=1.9cm,yshift=-12.85cm]current page.north west)
    {\sffamily\bfseries\fontsize{124}{124}\selectfont 9};
  \draw[milcGradeText,line width=1.7pt]
    ([xshift=4.52cm,yshift=-11.68cm]current page.north west) circle (.13cm);

  \node[anchor=west,text=milcGradeText,text width=8.7cm,align=left] at ([xshift=10.35cm,yshift=-11.6cm]current page.north west)
    {
     {\sffamily\bfseries\fontsize{19}{22}\selectfont Herramientas digitales ---\\Figma, Canva\\o Adobe Express}\\[4mm]
     {\sffamily\bfseries\fontsize{23}{26}\selectfont Guía 16-9-TIC}\\[2mm]
     {\sffamily\fontsize{13}{16}\selectfont Período 2 · Diseño editorial digital}\\[5mm]
     {\sffamily\bfseries\fontsize{11}{13}\selectfont Institución Educativa Sor María Juliana}\\[1mm]
     {\sffamily\fontsize{10}{12}\selectfont Autor: Dr. Álvaro Cárdenas Orozco}
    };

  \path[fill=milcGradeSoft,rounded corners=7mm]
    ([xshift=1.35cm,yshift=.85cm]current page.south west) rectangle
    ([xshift=-1.35cm,yshift=4.25cm]current page.south east);
  \draw[milcGradeDark,line width=.65pt,rounded corners=7mm]
    ([xshift=1.35cm,yshift=.85cm]current page.south west) rectangle
    ([xshift=-1.35cm,yshift=4.25cm]current page.south east);
  \node[anchor=center,text=milcNegro,text width=17.0cm,align=center] at ([yshift=2.55cm]current page.south)
    {\sffamily\fontsize{10}{12}\selectfont Metodología MILC · Apertura ancestral · Pensamiento computacional · Triángulo Dussel-Estoicismo-Floridi\\
    \textcolor{milcGradeDark}{\bfseries Producto:} Portada de la revista del periodo hecha en la herramienta digital elegida, usando grilla, tipografía, paleta y lógica imagen-texto de las sesiones anteriores};
\end{tikzpicture}
\mbox{}
\newpage

\section*{Apertura: Herramientas digitales --- Figma, Canva o Adobe Express}

\begin{softbox}{milcGradeDark}{milcGradeSoft}{Saber ancestral}
El paso del \textbf{taller físico al taller digital} no cambió el oficio: cambió la herramienta. Los \emph{impresores antiguos} del Valle también dominaban una herramienta compleja --- la prensa de tipos móviles, los rodillos, las planchas, las tintas. El \emph{tipógrafo de los años 50} usaba linotipia: una máquina que componía líneas enteras de plomo fundido. El oficio era duro, lento, costoso. Aprender a usar la herramienta requería años. Hoy Figma, Canva o Adobe Express son las nuevas linotipias: herramientas complejas que aceleran el oficio sin reemplazar el criterio. Saber usarlas con maestría es heredar el oficio editorial en formato 2026.
\end{softbox}

\begin{softbox}{milcTurquesa}{milcGris}{Saber contemporáneo}
Las herramientas más usadas hoy: \textbf{Figma} (gratis para individuos, profesional, colaborativo, el estándar de la industria), \textbf{Canva} (más sencilla, con plantillas, ideal para principiantes), \textbf{Adobe Express} (versión simplificada de Adobe, gratuita con cuenta). Cada una tiene fortalezas distintas. Lo importante es que la herramienta no reemplaza el criterio: si llegas a Figma sin haber pensado grilla, tipografía y paleta, vas a producir piezas confusas con software profesional.
\end{softbox}

\begin{softbox}{milcMostaza}{milcGris}{Pregunta-puente}
\textit{¿Qué herramienta digital ya usas (aunque sea Word, Google Docs, o el editor de Instagram)? ¿En qué te ayuda y qué decisiones todavía dependen de tu criterio?}
\end{softbox}

\begin{softbox}{milcVerde}{milcGris}{Saber y saber hacer}
Al terminar podrás: (1) \textbf{identificar} 3 herramientas digitales editoriales y sus diferencias; (2) \textbf{aplicar} grilla + tipografía + color en una herramienta real; (3) \textbf{crear} la portada de tu revista del periodo en la herramienta elegida; (4) \textbf{evaluar} con criterio si una herramienta sirve a tu propósito o te obliga a forzar el diseño.
\end{softbox}

\small
\begin{tabularx}{\textwidth}{>{\bfseries}p{3.0cm}Y}
\rowcolor{milcGradePrimary!12}Autor & Dr. Álvaro Cárdenas Orozco\\
DBA & Producción de piezas editoriales digitales con criterio estético y ético (MEN, T\&I)\\
Referentes & Tipografía colombiana · Diseño centrado en accesibilidad · Floridi\\
Producto final & Portada de la revista del periodo hecha en la herramienta digital elegida, usando grilla, tipografía, paleta y lógica imagen-texto de las sesiones anteriores\\
\end{tabularx}
\normalsize

\puente{Antes de empezar, mira el viaje completo. Las sesiones 1-5 te dieron criterio editorial. Hoy aplicas todo eso en una herramienta digital real. La sesión 7 será sobre accesibilidad y luego maquetar la revista completa.}

\begin{titlebox}{milcGradeDark}
Ruta MILC: así vamos a aprender
\end{titlebox}
\begin{tabularx}{\textwidth}{>{\centering\arraybackslash}X>{\centering\arraybackslash}X>{\centering\arraybackslash}X>{\centering\arraybackslash}X}
\phasecard{1}{milcVerde}{milcGris}{Escuta\\[-1mm]\small Observo, escucho y pregunto.} &
\phasecard{2}{milcTurquesa}{milcGris}{Sistematización\\[-1mm]\small Pensamiento computacional.} &
\phasecard{3}{milcMagenta}{milcGris}{Praxis\\[-1mm]\small Construyo y aplico.} &
\phasecard{4}{milcVino}{milcGris}{Evaluación\\[-1mm]\small Reflexión liberadora.}
\end{tabularx}


\puente{Antes de elegir tu herramienta, vamos a comparar las 3 más usadas. Cada una sirve a perfiles distintos: Figma para quien quiere profundidad profesional, Canva para velocidad con plantillas, Adobe Express como balance.}

\begin{titlebox}{milcVerde}
Fase 1 · Escuta
\end{titlebox}

\begin{softbox}{milcVerde}{milcGris}{Escena para escuchar}
\textbf{Actividad 1 · IDENTIFICA --- 3 herramientas digitales editoriales} (15 min · individual). Abre las páginas de \textbf{figma.com}, \textbf{canva.com} y \textbf{express.adobe.com}. En cada una: ¿qué muestra la página principal? ¿qué tipo de usuario apunta? ¿es gratis o paga? ¿qué tipo de pieza es el primer ejemplo? Crea cuenta gratuita en una de las 3 (la que más se acerque a tu necesidad).
\end{softbox}

\checkbox Visité las 3 páginas (Figma, Canva, Adobe Express).\\
\checkbox Comparé tipo de usuario, precios y ejemplos.\\
\checkbox Creé cuenta gratuita en una de las 3.

\begin{infoband}{milcVerde}
\textbf{Lo que va al cuaderno (Actividad 1 · IDENTIFICA).} Título: «Actividad 1 --- 3 herramientas editoriales digitales». \textbf{Formato:} tabla de 4 columnas (Herramienta~|~Tipo de usuario~|~Costo (gratis/paga)~|~Mi impresión), 3 filas. Al final indica cuál elegiste y por qué (1 frase). \textbf{Sabes que terminaste cuando} la elección está justificada según tu necesidad real, no por moda.
\end{infoband}

\puente{Ya viste las 3 herramientas. Ahora vamos a entender los \textbf{4 elementos básicos} que cualquier herramienta editorial digital ofrece: lienzo + capas + tipografía + exportación. Si dominas estos 4 en cualquier herramienta, dominas el oficio digital.}

\begin{titlebox}{milcTurquesa}
Fase 2 · Sistematización con pensamiento computacional
\end{titlebox}

\begin{softbox}{milcTurquesa}{milcGris}{Cuatro pilares aplicados al tema}
Toda herramienta editorial digital opera con 4 elementos básicos. Vamos a descomponerlos con los pilares del pensamiento computacional para que aprendas la herramienta más rápido.
\end{softbox}

\small
\begin{tabularx}{\textwidth}{>{\bfseries\RaggedRight\arraybackslash}p{3.6cm}Y}
\rowcolor{milcTurquesa!90}\color{white}Pilar & \color{white}\bfseries Aplicación al tema\\
1. Descomposición & Las herramientas digitales editoriales se descomponen en 4 elementos: (1) \textbf{lienzo} (la página o frame de trabajo, con tamaño definido), (2) \textbf{capas} (cada elemento es una capa apilable y editable independientemente), (3) \textbf{tipografía} (sistema para elegir, escalar y combinar fuentes), (4) \textbf{exportación} (formato final --- PNG, JPG, PDF, video).\\
2. Reconocimiento de patrones & Toda herramienta sigue el patrón \textbf{``modelo de capas''}: el texto, la imagen, los colores son objetos independientes que ocupan capas. Mover una capa no afecta a otra. Aprender a pensar en capas es la primera habilidad digital editorial.\\
3. Abstracción & Lo esencial cabe en una pregunta: \emph{¿cuál herramienta sirve a la pieza que necesito y a mi nivel actual?}. Sin esa pregunta, eliges por marca o por moda.\\
4. Algoritmo conceptual & Pasos para empezar en cualquier herramienta: (1) crea lienzo con tamaño correcto (A4, Instagram post, etc.); (2) agrega grilla guía; (3) agrega capas de texto con tu par tipográfico; (4) agrega capas de imagen (foto, ilustración); (5) aplica colores de tu paleta; (6) exporta en formato adecuado.\\
\end{tabularx}
\normalsize

\begin{drawbox}{milcTurquesa}{Anatomía de tu primer documento digital}
\textbf{Lienzo:} tamaño A4, A5, Instagram post (1080×1080), historia (1080×1920), etc. \\[1mm] \textbf{Grilla guía:} columnas y márgenes invisibles para alinear elementos. \\[1mm] \textbf{Capas:} cada texto, imagen y forma es una capa editable independiente. \\[1mm] \textbf{Tipografía:} aplica tu par display + cuerpo de Sesión 3. \\[1mm] \textbf{Color:} aplica tu paleta HEX de Sesión 4. \\[1mm] \textbf{Exportación:} JPG/PNG para web, PDF para imprimir.
\end{drawbox}

\begin{softbox}{milcMostaza}{milcGris}{Errores comunes y su corrección}
(1) Empezar a diseñar sin grilla guía: todo termina desalineado. (2) Usar tipografías y colores diferentes a los del sistema visual establecido (sesiones 3 y 4): la revista pierde coherencia. (3) Aplastar todo en una sola capa: imposible editar luego. (4) Exportar en formato equivocado (PDF cuando necesitas JPG, o al revés). (5) Llenar con stickers y elementos decorativos que distraen del contenido editorial.
\end{softbox}

\begin{infoband}{milcTurquesa}
\textbf{Lo que va al cuaderno (Actividad 2 · ANALIZA).} Título: «Actividad 2 --- 4 elementos de una herramienta digital editorial». \textbf{Formato:} 4 fichas, una por elemento (lienzo, capas, tipografía, exportación), cada una con: definición + cómo lo aplica la herramienta que elegiste + captura de pantalla o dibujo del menú. \textbf{Sabes que terminaste cuando} puedes ubicar los 4 elementos en la herramienta sin titubear.
\end{infoband}

\puente{Tienes el método. Toca producir: la portada de tu revista del periodo en la herramienta que elijas. Aplicando grilla, tipografía, color y relación imagen-texto. Es el primer producto en software de tu revista.}

\begin{titlebox}{milcMagenta}
Fase 3 · Praxis con pensamiento computacional
\end{titlebox}

\begin{softbox}{milcMagenta}{milcGris}{Construyendo el producto con los cuatro pilares}
\textbf{Actividad 3 · CREA --- Portada digital de tu revista} (35 min · individual). Hora de producir en software. Vas a diseñar la portada de tu revista del periodo en la herramienta digital elegida, aplicando grilla + tipografía + paleta + relación imagen-texto.
\end{softbox}

\small
\begin{tabularx}{\textwidth}{>{\bfseries\RaggedRight\arraybackslash}p{3.6cm}Y}
\rowcolor{milcMagenta!90}\color{white}Pilar & \color{white}\bfseries Aplicación al producto\\
Descomposición & Tu portada se descompone en 5 elementos: (1) lienzo A4 vertical, (2) grilla guía, (3) imagen principal (foto, ilustración o composición visual), (4) titular principal con tu display tipográfica, (5) datos de pie (nombre de revista, número, fecha).\\
Patrones & Sigue el patrón \textbf{``aplica el sistema antes que improvisa''}: antes de tocar cualquier elemento, configura tu par tipográfico y tu paleta como estilos en la herramienta. Eso garantiza coherencia.\\
Abstracción & Cada elemento expresa una sola idea. Si tu portada quiere decir 3 cosas distintas, divídela en 3 portadas o elige una.\\
Algoritmo de armado & Algoritmo: (1) crea lienzo A4 vertical en la herramienta; (2) agrega grilla guía con columnas (idealmente 3); (3) ubica imagen principal (foto subida o ilustración elegida); (4) escribe titular con tu display tipográfica; (5) aplica tu paleta de color; (6) agrega datos de pie (nombre de revista + número/fecha); (7) exporta como JPG o PDF.\\
\end{tabularx}
\normalsize

\begin{drawbox}{milcMagenta}{Checklist antes de exportar la portada}
\checkbox Lienzo A4 vertical (210×297mm) o tamaño revista coherente. \\[1mm] \checkbox Grilla guía visible mientras diseñas (puedes ocultarla al exportar). \\[1mm] \checkbox Tipografía display de Sesión 3 aplicada al titular. \\[1mm] \checkbox Paleta de color de Sesión 4 aplicada (mínimo dominante + acento). \\[1mm] \checkbox Imagen principal compositivamente fuerte (no relleno). \\[1mm] \checkbox Datos de pie con nombre de revista + número/fecha. \\[1mm] \checkbox Exportada en formato adecuado (JPG/PDF).
\end{drawbox}

\begin{softbox}{milcMagenta}{milcGris}{Plantilla de trabajo}
\textbf{Herramienta usada.} \textit{Figma / Canva / Adobe Express.} \\[1mm] \textbf{Tamaño del lienzo.} \textit{A4 vertical o el que elegí.} \\[1mm] \textbf{Imagen principal.} \textit{Descripción breve.} \\[1mm] \textbf{Titular.} \textit{Texto exacto.} \\[1mm] \textbf{Aplicación del sistema visual.} \textit{Cómo apliqué grilla, tipografía y color.}
\end{softbox}

\begin{infoband}{milcMagenta}
\textbf{Lo que va al cuaderno (Actividad 3 · CREA).} Título: «Actividad 3 --- Mi portada digital de revista». \textbf{Formato:} captura impresa o dibujada de la portada + 1 párrafo de proceso (qué fue difícil, qué fue fluido). \textbf{Extensión:} 1 página. \textbf{Sabes que terminaste cuando} la portada se ve coherente con tus decisiones editoriales previas (no es ``otra cosa'').
\end{infoband}

\puente{Lo que sigue es el resumen de qué entregas hoy: portada digital. El criterio principal: que aplique conscientemente las decisiones editoriales previas, no que sea ``vistosa''.}

\begin{titlebox}{milcMostaza}
Producto de la sesión
\end{titlebox}
\textbf{Portada digital de revista hecha en herramienta editorial con sistema visual aplicado}.
\begin{softbox}{milcMostaza}{milcGris}{Criterios de calidad}
(1) Portada producida en software (Figma/Canva/Adobe Express). (2) Tamaño A4 vertical o equivalente. (3) Tipografía display de Sesión 3 aplicada. (4) Paleta de color de Sesión 4 aplicada. (5) Imagen principal y titular dialogan. (6) Datos de pie completos. (7) Exportada en formato adecuado para compartir/imprimir.
\end{softbox}

\puente{Antes de cerrar, mira las herramientas digitales desde las cinco dimensiones humanas. Cada herramienta es un sistema económico (Adobe paga, Figma gratis condicional), político (quién tiene acceso) y cognitivo (qué facilita y qué dificulta).}

\begin{titlebox}{milcVino}
Fase 4 · Evaluación liberadora — cinco dimensiones
\end{titlebox}

\begin{softbox}{milcVino}{milcGris}{Sentido de la evaluación}
Evaluar no es solo revisar una nota. Es reconocer qué aprendí, qué cambió en mí, qué emociones manejé, cómo me relaciono con la ciudad y la comunidad, y qué le dejaré a quien viene detrás.
\end{softbox}

\textbf{Mi reflexión liberadora en cinco dimensiones.} Completa cada frase iniciada:

\small
\begin{tabularx}{\textwidth}{>{\bfseries\RaggedRight\arraybackslash}p{3.4cm}Y>{\RaggedRight\arraybackslash}p{4.6cm}}
\rowcolor{milcVino}\color{white}Dimensión & \color{white}\bfseries Mi respuesta (frame starter) & \color{white}\bfseries Algo que descubrí\\
1. Personal & Lo que más cambió en mí fue \linead{6.5cm} & \linead{4.2cm}\\
2. Emocional & La emoción más fuerte fue \linead{2.5cm} y la regulé \linead{3cm} & \linead{4.2cm}\\
3. Ciudadana & En democracia esto importa porque \linead{6cm} & \linead{4.2cm}\\
4. Local (Cartago) & En mi barrio sería útil para \linead{6.5cm} & \linead{4.2cm}\\
5. Intergeneracional & A mi abuela/abuelo le diría \linead{6.5cm} & \linead{4.2cm}\\
\end{tabularx}
\normalsize

\begin{infoband}{milcVino}
\textbf{Para la próxima sustentación.} Vuelve a esta tabla en seis meses. Verás que las respuestas cambian — no porque lo aprendido cambie, sino porque tú cambias. La evaluación liberadora es una conversación contigo a través del tiempo.
\end{infoband}


\puente{Y para terminar, tres voces sobre herramientas digitales. Dussel sobre quién tiene acceso a las profesionales, Marco Aurelio sobre la herramienta que no debe gobernarte, Floridi sobre las plataformas como infraestructura.}

\begin{titlebox}{milcGradeDark}
Triángulo de pensamiento · cierre filosófico
\end{titlebox}

\begin{softbox}{milcVino}{milcGris}{Dussel · lente del nosotros}
\textit{``El acceso a las herramientas profesionales es la nueva alfabetización del XXI; quien no las domina depende del que las domina.''}

Figma es gratis para individuos. Canva tiene plan gratuito robusto. Adobe Express también. Eso significa que la barrera de acceso a las herramientas profesionales de diseño editorial es \emph{conexión + tiempo de aprendizaje}, no dinero. Quien tiene acceso a internet y disciplina puede llegar al mismo nivel técnico que un estudiante de diseño en Bogotá o Nueva York. Aprovechar esa ventana es soberanía técnica.

\textbf{¿Estoy aprovechando la ventana actual de acceso gratuito a herramientas profesionales? ¿O me limito a las herramientas básicas porque ``son más fáciles''?}
\end{softbox}
\linead{\linewidth}\\[1mm]
\linead{\linewidth}

\begin{softbox}{milcMostaza}{milcGris}{Estoicismo · Marco Aurelio · cuidado interior}
\textit{``Aprende a usar la herramienta sin que ella te use a ti.''}

Cada herramienta tiene su lógica: Canva te empuja a plantillas, Figma te empuja a sistemas complejos, Adobe te empuja a su ecosistema. La sabiduría estoica del diseñador: \emph{conocer el sesgo de tu herramienta} para usarlo o resistirlo conscientemente. Si la herramienta tiene 1000 plantillas, la pregunta es si las usas o las dejas decidir por ti.

\textbf{Cuando uso una herramienta digital, ¿estoy decidiendo yo o estoy siguiendo lo que la herramienta me sugiere? ¿Dónde está el criterio?}
\end{softbox}
\linead{\linewidth}\\[1mm]
\linead{\linewidth}

\begin{softbox}{milcTurquesa}{milcGris}{Floridi · lente de la infoesfera}
\textit{``Las plataformas digitales no son neutrales: cada una organiza la creación según su propia lógica económica y estética.''}

Figma se diseñó para colaboración en empresas tech. Canva se diseñó para usuarios sin experiencia que necesitan resultados rápidos. Adobe se diseñó para profesionales del mercado creativo. Tus piezas saldrán parecidas a las de quien usa la misma plataforma, porque las plataformas \emph{sesgan la creación}. Reconocer ese sesgo es ciudadanía digital avanzada.

\textbf{¿Mis trabajos en Canva se parecen a otros trabajos en Canva? ¿Y mis trabajos en Figma a otros en Figma? ¿Cuál es mi voz propia más allá de la herramienta?}
\end{softbox}
\linead{\linewidth}\\[1mm]
\linead{\linewidth}

\begin{titlebox}{milcVino}
Compromiso final
\end{titlebox}

\small
\begin{tabularx}{\textwidth}{>{\RaggedRight\arraybackslash}p{3.0cm}YYY}
\rowcolor{milcVino}\color{white}\bfseries Criterio & \color{white}\bfseries Lo logré & \color{white}\bfseries En proceso & \color{white}\bfseries Necesito apoyo\\
Comprensión del tema & Explico ideas centrales con mis palabras. & Reconozco ideas pero falta precisión. & Necesito repasar conceptos base.\\
Pensamiento computacional & Apliqué los 4 pilares al diseñar mi producto. & Apliqué algunos pilares. & Necesito releer la fase 2.\\
Aplicación y producto & Mi producto cumple los criterios y se puede revisar. & Producto funcional con ajustes pendientes. & Aún en construcción.\\
Reflexión 5 dimensiones & Respondí cada dimensión con honestidad. & Respondí algunas con detalle. & Mis respuestas son muy generales.\\
Triángulo de pensamiento & Conecté las 3 voces con mi trabajo real. & Conecté al menos una. & No relaciono las voces con mi proceso.\\
\end{tabularx}
\normalsize

\textbf{Hoy entendí que:}\\[1mm]
\linead{\linewidth}\\[1mm]
\linead{\linewidth}

\textbf{Mi compromiso para la próxima sesión es:}\\[1mm]
\linead{\linewidth}\\[1mm]
\linead{\linewidth}

\begin{softbox}{milcMostaza}{milcGris}{Cita de despedida · Marco Aurelio}
\textit{``El obstáculo en el camino se vuelve el camino. Lo que estorba al camino se vuelve camino.''}\\[1mm]
Cada error de hoy, bien observado, se convierte en pieza del camino que pisarás mañana.
\end{softbox}

\small
\begin{tabularx}{\textwidth}{>{\bfseries}p{3.6cm}Y}
\rowcolor{milcGradeDark!12}Nombre completo: & \linead{10cm}\\
Fecha: & \linead{4cm} \quad Curso/grupo: \linead{4cm}\\
Mi firma: & \linead{9cm}\\
\end{tabularx}
\normalsize

\begin{infoband}{milcGradeDark}
\textbf{Para la próxima semana.} Hace una segunda versión de tu portada cambiando UN solo elemento (color principal, tipografía, posición de la imagen, titular). Compara las dos versiones. Anota cuál se siente mejor y por qué. Ese ejercicio de iteración es lo que separa al diseñador del usuario de plantillas. El próximo lunes traemos las dos versiones al aula.
\end{infoband}

\end{document}
